\subsection{全对称斜晶格与多方向增益}

\begin{proof}[命题~\ref{prop:nonexact} 的证明]
固定 \(n\ge2,\ \kappa>0\)，设
\[
t=u^n,\quad \delta=\frac{u}{1+2u},\quad r_0=\frac1{1+2u}=1-2\delta,
\quad a=(2u)^n,\quad v=2u,\quad e=\kappa u^4.
\]
令 \(J\) 为全 \(1\) 矩阵，
\[
T=(1-e)I+eJ,\qquad
\Lambda=T\mathbb Z^n,\qquad
d=\det T=(1-e)^{n-1}(1+(n-1)e),\qquad
\mu=(2\delta)^n=a r_0^n.
\]
所有孔仍为轴平行立方体：
\[
D=\bigcup_{\lambda\in\Lambda}(\lambda+[-\delta,\delta]^n),
\qquad E_\infty=\mathbb R^n\setminus D.
\]
这是全对称斜晶格，不是单位立方体铺砌。

取 \(M=3^n\)。以下显式条件足够：
\begin{equation}\label{eq:skew-1-1}
u\le\frac14,\qquad
(n-1)\kappa u^3<\frac13,\qquad
M\kappa^2u^5<\frac{4n}{9}.
\end{equation}
它们给出
\begin{equation}\label{eq:skew-1-2}
e<\delta/2,\qquad
(n-1)e<1/12,\qquad r_0\ge2/3>2\delta,\qquad
d-\mu>11/12-1/9=29/36.
\end{equation}
因此孔内部互不交，所有后续分母为正。

定义
\[
K=\frac{a}{(2\delta)^{n-2}}=4u^2r_0^{2-n},\qquad
F(h)=K[(r_0+h)^n-r_0^n].
\]
记 \(h_*,h_M\) 为
\[
F(h_*)=e^2,\qquad F(h_M)=Me^2,
\]
即
\[
h_c=(r_0^n+ce^2/K)^{1/n}-r_0,\qquad c\in\{1,M\}.
\]
再设
\[
A=2\int_0^{h_*}\left[e-\frac{F(h)}e\right]dh
\]
和单个无序坐标对的价格
\begin{equation}\label{eq:skew-1-3}
B_{\rm pair}=2A(2\delta)^{n-2}.
\end{equation}
总失败量上界为
\begin{equation}\label{eq:skew-1-4}
B_n=\binom n2 B_{\rm pair}
+\begin{cases}
0,&n=2,\\
8\binom n3(e+h_M)^3(2\delta)^{n-3},&n\ge3.
\end{cases}
\end{equation}

先证明以下可实现的定量下界： 在 \eqref{eq:skew-1-1} 下，
\begin{equation}\label{eq:skew-1-5}
 C_n(1-t)\ge\frac{d-B_n}{d-\mu}.
\end{equation}
固定 \(n,\kappa\) 后，
\begin{equation}\label{eq:skew-1-6}
\liminf_{t\downarrow0}\frac{C_n(1-t)-L_n(t)}{av^8}
\ge\binom n2\,
\frac{\kappa^2-\kappa^3/(8n)}{256}.
\end{equation}
\paragraph{穷尽可能进入见证的所有晶格孔}

对中心孔内 \(x\)，设
\[
h_{i,\sigma}=\delta-\sigma x_i,\quad
s=\min_{i,\sigma}h_{i,\sigma}=\delta-\|x\|_\infty,\quad
r=r_0+s,\quad \sigma\in\{+1,-1\}.
\]
于是 \(r\ge r_0>2\delta\)、\(r+\delta=1-\|x\|_\infty\)。

\paragraph{大整数指数全部排除}

\(\|T-I\|_{\infty\to\infty}=(n-1)e<1/12\)，所以
\[
\|Tk\|_\infty\ge[1-(n-1)e]\|k\|_\infty.
\]
当 \(\|k\|_\infty\ge2\) 时，
\[
\|Tk-x\|_\infty
\ge2[1-(n-1)e]-\|x\|_\infty
>1-\|x\|_\infty=r+\delta.
\]
因此对应孔与见证不相交。孔内部互不交也由同一分离下界及 \(2\delta<1-(n-1)e\) 得到。

\paragraph{小整数指数中只有平衡指数可以进入}

现在 \(k\in\{-1,0,1\}^n\)，令 \(S=\sum_i k_i\)。如果 \(S>0\)，存在
\(k_i=+1\)，该坐标给
\[
(Tk)_i=(1-e)+eS=1+e(S-1)\ge1.
\]
如果 \(S<0\)，存在 \(k_i=-1\)，并有
\[
-(Tk)_i=1+e(|S|-1)\ge1.
\]
故任意 \(S\ne0\) 的指数都在某坐标至多边界接触，不增加正体积。

若 \(S=0\)，则 \(Tk=(1-e)k\)。非零坐标数
\(q=q(k)\) 必为偶数，且正、负坐标各 \(q/2\) 个。定义
\begin{equation}\label{eq:skew-2-1}
L_{i,\sigma}=(e+s-h_{i,\sigma})_+.
\end{equation}
每个 \(L_{i,\sigma}\le e<2\delta\)。对应孔在非零坐标的精确交长为
\(L_{i,k_i}\)，零坐标上则完整包含宽 \(2\delta\)。所以总交量精确为
\[
|D\cap Q(x,r)|=\mu+(2\delta)^{n-2}R(x),
\]
\begin{equation}\label{eq:skew-2-2}
R(x)=
\sum_{\substack{k\in\{-1,0,1\}^n\setminus\{0\}\\ \sum k_i=0}}
(2\delta)^{2-q(k)}
\prod_{i:k_i\ne0}L_{i,k_i}.
\end{equation}
这没有丢掉 \(q\ge4\) 的高阶平衡邻孔。实际严格见证条件等价于
\begin{equation}\label{eq:skew-2-3}
R(x)<F(s).
\end{equation}

\paragraph{先约束整个坏集，排除同坐标的相反符号}

由于 \(e<2\delta\)，\eqref{eq:skew-2-2} 中每一项至多
\[
e^2\left(\frac e{2\delta}\right)^{q(k)-2}\le e^2.
\]
项数小于 \(3^n=M\)，从而 \(R(x)\le Me^2\)。在孔内部 \(s>0\)，若见证失败，则
\begin{equation}\label{eq:skew-3-1}
0<F(s)\le R(x)\le Me^2,\qquad s\le h_M.
\end{equation}

守卫 \eqref{eq:skew-1-1} 还给
\begin{equation}\label{eq:skew-3-2}
 F(\delta-e)>2Me^2,\qquad h_M<\delta-e.
\end{equation}
事实上 \(F'(h)\ge4nu^2r_0\)，而 \(e<\delta/2\)，所以
\[
F(\delta-e)>2nu^2r_0\delta
=2nu^3r_0^2
\ge\frac{8n}{9}u^3>2Me^2.
\]
最后一步是 \(M\kappa^2u^5<4n/9\)。

在坏集上，如果 \(L_{i,\sigma}>0\)，则
\begin{equation}\label{eq:skew-3-3}
h_{i,\sigma}<e+s\le e+h_M<\delta.
\end{equation}
同一坐标的两侧缺口之和为 \(2\delta\)，故 \emph{在坏集上每坐标至多一个符号的 \(L\) 为正}。这直接排除了“同一个无序坐标对的两个相反短角同时正，却只涉及两个坐标”的遗漏。该断言仅在坏集局部化之后使用，没有把它当作所有点的几何事实。

\paragraph{活跃短角的穷尽二分}

把有向短角 \((i,j)\)、\(i\ne j\)，称为活跃，当且仅当
\[
L_{i,+}>0,\qquad L_{j,-}>0.
\]
它对应 \(q=2\) 平衡指数 \(k=e_i-e_j\)。

任何正的 \(q\ge4\) 平衡项都有至少两个正坐标和至少两个负坐标。其所有相关交长均正，因此任取一个正坐标和一个负坐标都会产生一个活跃 \(q=2\) 短角；至少有四条不同活跃短角。于是：

\begin{itemize}
\item 坏点至少有一个活跃短角，否则 \eqref{eq:skew-2-2} 的所有项都为零，与 \eqref{eq:skew-3-1} 矛盾。
\item 若恰有一个活跃短角，所有 \(q\ge4\) 项都为零，总额外归一化质量只有该角的 \(L_{i,+}L_{j,-}\)。
\item 若至少有两个不同活跃短角，则它们至少涉及三个不同坐标。若只涉及同两个坐标，两角只能方向相反，违背 \eqref{eq:skew-3-3} 的单坐标符号互斥。
\end{itemize}

这三点同时证明了高阶平衡邻孔没有脱离短角计数的独立坏机制。

\paragraph{唯一活跃短角：单角面积价格}

固定唯一活跃角 \((i,j)\)，令
\(\alpha=h_{i,+}\)、\(\beta=h_{j,-}\)。由于坏点仅含该贡献，
\[
F(s)\le L_{i,+}L_{j,-}\le e^2,\qquad
s\le h_*,\qquad \alpha,\beta<e+h_*<\delta.
\]
若最小面来自这两个坐标，在 \(\alpha\le\beta,s=\alpha\) 分支，
\[
L_{i,+}L_{j,-}=e(e+\alpha-\beta),\qquad
\alpha\le\beta\le\alpha+e-F(\alpha)/e.
\]
交换两坐标，得到同一有向角的二维面积上界
\(A=2\int_0^{h_*}[e-F(h)/e]dh\)。其他坐标最多乘
\((2\delta)^{n-2}\)。

最小面不可能来自其他坐标。否则对某 \(k\notin\{i,j\}\) 和符号
\(\sigma\) 有 \(h_{k,\sigma}=s\)，从而 \(L_{k,\sigma}=e>0\)。
若 \(\sigma=+\)，则 \((k,j)\) 也是活跃短角；若 \(\sigma=-\)，
则 \((i,k)\) 也是活跃短角。两种情形都与唯一性矛盾。
因此单角坏点只需前面两个关键面的面积价格，不再加入额外维度面层项。
每个无序坐标对有两个方向，合计为 \(B_{\rm pair}=2A(2\delta)^{n-2}\)。
对 \(\binom n2\) 个无序对求和，得到 \eqref{eq:skew-1-4} 的第一项。

\paragraph{多个活跃短角：至少三个面层}

由前述二分，多个活跃角的坏点必有三个不同坐标
\(i_1,i_2,i_3\)，各有一个指定符号 \(\sigma_\ell\)，满足
\[
0\le h_{i_\ell,\sigma_\ell}<e+h_M.
\]
固定三坐标和符号，该面层交集体积至多
\((e+h_M)^3(2\delta)^{n-3}\)。
遍历 \(\binom n3\) 个坐标三元组及 \(2^3\) 个符号模式，得到
\begin{equation}\label{eq:skew-4-1}
|\{\text{多个活跃角的坏点}\}|
\le8\binom n3(e+h_M)^3(2\delta)^{n-3}.
\end{equation}
这是允许重叠的联合上界，无需选择唯一三元组或证明各层互不交。
当 \(n=2\) 时多个活跃角坏点不存在，\eqref{eq:skew-4-1} 项不出现。

结合两类及孔边界的零测集，所选半径失败的总体积至多 \(B_n\)。

\paragraph{有限测度输入、严格水平及定量结论}

若 \(B_n\ge \mu\)，\eqref{eq:skew-1-5} 右端不超过 \(1\)，由 \(C_n\ge1\) 直接成立。以下设 \(B_n<\mu\)。

令
\[
\Omega_N=T([-N,N]^n),\qquad E_N=E_\infty\cap\Omega_N,\qquad N\in\mathbb N,\quad N\ge4.
\]
这保留所有被边界切开的周期孔。变换后
\(\mathbf1_{E_\infty}(Ty)\) 在每个坐标以 \(1\) 为周期，每个基本胞的缺陷体积为 \(\mu\)，所以
\begin{equation}\label{eq:skew-5-1}
|E_N|=(2N)^n(d-\mu)
\end{equation}
精确成立。它不是“只挖去内部索引孔”的另一输入。

由 \(\|T-I\|_{\infty\to\infty}<1/12\)，Neumann 级数给
\[
\|T^{-1}\|_{\infty\to\infty}
\le[1-(n-1)e]^{-1}<2.
\]
对 \(\|k\|_\infty\le N-3\) 的孔，任意孔点 \(x=Tk+y\) 的所选见证均位于 \(\Omega_N\)，因为
\[
\|T^{-1}z-k\|_\infty
\le\|T^{-1}\|_{\infty\to\infty}(\delta+r(x))<2
\]
对见证内每个 \(z\) 成立。所有好点原来的严格平均不等式完整保留。
内部孔至少有 \((2N-5)^n\) 个；密度定理又给 \(E_N\) 几乎处处位于自身严格 halo。故
\begin{equation}\label{eq:skew-5-2}
\frac{|\Hc(1-t;E_N)|}{|E_N|}
\ge1+\frac{(2N-5)^n(\mu-B_n)}{(2N)^n(d-\mu)}
\ge\frac{d-B_n}{d-\mu}-O_n(\mu/N).
\end{equation}
令 \(N\to\infty\) 得 \eqref{eq:skew-1-5}；选择 \(N=\lceil u^{-9}\rceil\) 则
\(O_n(\mu/N)=o(av^8)\)。

好点集由严格不等式 \eqref{eq:skew-2-3} 定义，可测且使用有限真实半径。等号点可全部计为坏点；孔边界 \(s=0\) 为零测集。上述证明从未将非严格平均升级为严格见证，也不需要所有好点具有统一平均裕量。

\paragraph{一般维数渐近}

固定 \(n,\kappa\)，由根方程得到：
\[
h_*\sim\frac{e^2}{4nu^2r_0},\qquad
h_M\sim\frac{Me^2}{4nu^2r_0}=O_{n,\kappa}(u^6),\qquad
A\sim\frac{e^3}{4nu^2r_0}.
\]
因此
\[
\frac{B_{\rm pair}}{av^8}
\longrightarrow\frac{\kappa^3}{2048n}.
\]
新增联合项满足
\begin{equation}\label{eq:skew-6-1}
8\binom n3(e+h_M)^3(2\delta)^{n-3}
=O_{n,\kappa}(u^{n+9})=o(av^8).
\end{equation}
它比目标尺度多一个 \(u\) 幂，足以消失；不将它误写成精确坏量。

由特征值公式，
\[
1-d=\binom n2e^2+O_n(e^3).
\]
精确相减给
\[
\frac{d-B_n}{d-\mu}-L_n(t)
=\frac{\mu(1-d)-B_n(1-\mu)}{(d-\mu)(1-\mu)}.
\]
分母趋于 \(1\)，且
\[
\frac{\mu(1-d)}{av^8}\longrightarrow
\binom n2\frac{\kappa^2}{256}.
\]
结合 \eqref{eq:skew-5-2}、\eqref{eq:skew-6-1} 得 \eqref{eq:skew-1-6}。对 \(\kappa>0\) 优化右端多项式，最大值在 \(\kappa=16n/3\) 取得，
等于 \(\binom n2 n^2/27\)，得到命题~\ref{prop:nonexact}。
\end{proof}
